\subsection{\texorpdfstring{DEFINITION OF \(\mathfrak{sl}_2\) -TRIPLETS}{DEFINITION OF \textbackslash mathfrak\{sl\}\_2 -TRIPLETS}}

DEFINITION 1. An \(\mathfrak{sl}_2\) -triplet in \(\mathfrak{g}\) is a sequence \((x, h, y)\) of elements of \(\mathfrak{g}\) , distinct from \((0, 0, 0)\) , such that

\[
[ h, x ] = 2 x, [ h, y ] = - 2 y, [ x, y ] = - h.
\]

Let \((x,h,y)\) be an \(\mathfrak{sl}_{2}\) -triplet in g. The linear map \(\tau\) from \(\mathfrak{sl}(2,k)\) to g such that \(\tau(X_{+}) = x, \tau(H) = h, \tau(X_{-}) = y\) is a homomorphism which is non-zero and hence injective (since \(\mathfrak{sl}(2,k)\) is simple), and with image \(kx + kh + ky\) . We thus obtain a canonical bijection from the set of \(\mathfrak{sl}_{2}\) -triplets in g to the set of injective homomorphisms from \(\mathfrak{sl}(2,k)\) to g. If g is semi-simple and if \((x,h,y)\) is an \(\mathfrak{sl}_{2}\) -triplet in g, then x and y are nilpotent elements of g and h is a semi-simple element of g (Chap. I, §6, no. 3, Prop. 4).

Lemma 1. Let \(x, h, y, y' \in \mathfrak{g}\) . If \((x, h, y)\) and \((x, h, y')\) are \(\mathfrak{sl}_2\) -triplets in \(\mathfrak{g}\) , then \(y = y'\) .

Indeed, \(y - y' \in \mathrm{Ker}(\mathrm{ad}_{\mathfrak{g}}x)\) and \((\mathrm{ad}_{\mathfrak{g}}h)(y - y') = -2(y - y')\) . But \(\mathrm{ad}_{\mathfrak{g}}x\) is injective on \(\mathrm{Ker}(p + \mathrm{ad}_{\mathfrak{g}}h)\) for every integer \(p > 0\) (§1, no. 2, Cor. of Prop. 2).

Lemma 2. Let \(\mathfrak{n}\) be a subalgebra of \(\mathfrak{g}\) such that, for all \(n\in \mathfrak{n}\) , \(\operatorname{ad}_{\mathfrak{g}}(n)\) is nilpotent. Let \(h\in \mathfrak{g}\) be such that \([h,\mathfrak{n}] = \mathfrak{n}\) . Then \(e^{\operatorname{ad}_{\mathfrak{g}}\mathfrak{n}}.h = h + \mathfrak{n}\) .

It is clear that \(e^{\operatorname{ad}_{\mathfrak{g}}(\mathfrak{n})}.h \subset h + \mathfrak{n}\) . We shall prove that, if \(v \in n\) , then \(h + v \in e^{\operatorname{ad}_{\mathfrak{g}}(\mathfrak{n})}.h\) . It suffices to prove that \(h + v \in e^{\operatorname{ad}_{\mathfrak{g}}(\mathfrak{n})}.h + \mathscr {C} ^ {p} \mathfrak {n}\) for all \(p \geq 1\) (since \(C^{p}n = 0\) for sufficiently large p). This is clear for p = 1 since \(C^{1}n = n\) . Assume now that we have proved the existence of \(y_{p} \in n\) and \(z_{p} \in C^{p}n\) such that \(h + v = e^{\operatorname{ad}_{\mathfrak{g}}y_{p}}.h + z_{p}\) . Since \((\operatorname{ad}_{\mathfrak{g}}h)(\mathfrak{n}) = \mathfrak{n}\) , \((\operatorname{ad}_{\mathfrak{g}}h)|\mathfrak{n}\) is a bijection from n to n, hence its restriction to \(C^{p}n\) , which leaves \(C^{p}n\) stable, is also bijective; consequently, there exists \(z \in C^{p}n\) such that \(z_{p} = [z, h]\) . Then

\[
e ^ {\operatorname{ad} _ {\mathfrak {g}} \left(y _ {p} + z\right)} h - e ^ {\operatorname{ad} _ {\mathfrak {g}} y _ {p}} h \in [ z, h ] + \mathscr {C} ^ {p + 1} \mathfrak {n}
\]

SO

\[
e ^ {\operatorname{ad} _ {\mathfrak {g}} \left(y _ {p} + z\right)} h \in h + v - z _ {p} + [ z, h ] + \mathscr {C} ^ {p + 1} \mathfrak {n} = h + v + \mathscr {C} ^ {p + 1} \mathfrak {n}
\]

which establishes our assertion by induction on p.

Lemma 3. Let \(x \in \mathfrak{g}\) , \(\mathfrak{p} = \operatorname{Ker}(\operatorname{ad}x)\) , \(\mathfrak{q} = \operatorname{Im}(\operatorname{ad}x)\) . Then \([\mathfrak{p}, \mathfrak{q}] \subset \mathfrak{q}\) , and \(\mathfrak{p} \cap \mathfrak{q}\) is a subalgebra of \(\mathfrak{g}\) .

If \(u \in \mathfrak{p}\) and \(v \in \mathfrak{q}\) , there exists \(w \in \mathfrak{g}\) such that \(v = [x, w]\) , so

\[
[ u, v ] = [ u, [ x, w ] ] = [ x, [ u, w ] ] - [ [ x, u ], w ] = [ x, [ u, w ] ] \in \mathfrak {q}.
\]

On the other hand, p is a subalgebra of g, so \([p \cap q, p \cap q] \subset p \cap q\) .

Lemma 4. Let \((x,h,y)\) and \((x,h',y')\) be \(\mathfrak{sl}_2\) -triplets in \(\mathfrak{g}\) . There exists \(z\in \mathfrak{g}\) such that \(\mathrm{ad}_{\mathfrak{g}}z\) is nilpotent and such that

\[
e ^ {\mathrm{ad} _ {\mathfrak {g}} z} x = x, \quad e ^ {\mathrm{ad} _ {\mathfrak {g}} z} h = h ^ {\prime}, \quad e ^ {\mathrm{ad} _ {\mathfrak {g}} z} y = y ^ {\prime}.
\]

Let \(\mathfrak{n} = \mathrm{Ker}(\mathrm{ad} x) \cap \mathrm{Im}(\mathrm{ad} x)\) . For all \(p \in \mathbf{Z}\) , let \(\mathfrak{g}_p = \mathrm{Ker}(\mathrm{ad} h - p)\) . By §1, no. 3 (applied to the adjoint representation of \(kx + ky + kh\) on \(\mathfrak{g}\) ), we have that \(\mathfrak{n} \subset \sum_{p > 0} \mathfrak{g}_p\) , so \(\mathrm{ad}_{\mathfrak{g}} n\) is nilpotent for all \(n \in \mathfrak{n}\) , and \([h, \mathfrak{n}] = \mathfrak{n}\) . We have \([x, h' - h] = 0\) and \([x, y - y'] = h' - h\) , so \(h' - h \in \mathfrak{n}\) . By Lemmas 2 and 3, there exists \(z \in \mathfrak{n}\) such that \(e^{\mathrm{ad}_{\mathfrak{g}} z} h = h'\) . Since \(z \in \mathrm{Ker} \mathrm{ad}_{\mathfrak{g}} x\) , we have \(e^{\mathrm{ad}_{\mathfrak{g}} z} x = x\) . Lemma 1 now proves that \(e^{\mathrm{ad}_{\mathfrak{g}} z} y = y'\) . Q.E.D.

Let G be a group of automorphisms of g. Then two \(sl_{2}\) -triplets \((x,h,y)\) , \((x',h',y')\) are said to be G-conjugate if there exists \(g\in G\) such that \(gx=x'\) , \(gh=h',gy=y'\) .

PROPOSITION 1. Let G be a group of automorphisms of g containing \(\mathrm{Aut}_e(\mathfrak{g})\) . Let \((x,h,y)\) and \((x',h',y')\) be \(\mathfrak{sl}_2\) -triplets in g. Let

\[
\mathfrak {t} = k x + k h + k y, \quad \mathfrak {t} ^ {\prime} = k x ^ {\prime} + k h ^ {\prime} + k y ^ {\prime}.
\]

Consider the following conditions:

\begin{enumerate}
\def\labelenumi{(\roman{enumi})}
\item
  x and \(x'\) are G-conjugate;
\item
  \((x, h, y)\) and \((x', h', y')\) are G-conjugate;
\item
  \(\mathfrak{t}\) and \(\mathfrak{t}'\) are G-conjugate.
\end{enumerate}

We have (i) \(\Longleftrightarrow\) (ii) \(\Longrightarrow\) (iii). If \(k\) is algebraically closed, the three conditions are equivalent.

\begin{enumerate}
\def\labelenumi{(\roman{enumi})}
\item
  \(\Longleftrightarrow\) (ii): This follows from Lemma 4.
\item
  \(\Longrightarrow\) (iii): This is clear.
\end{enumerate}

We assume that \(k\) is algebraically closed and prove that (iii) \(\Longrightarrow\) (i). We treat first the case in which \(\mathfrak{t} = \mathfrak{t}' = \mathfrak{g} = \mathfrak{sl}(2, k)\) . Since \(\operatorname{ad}_{\mathfrak{g}} x\) is nilpotent, the endomorphism \(x\) of \(k^2\) is nilpotent (Chap. I, §6, Th. 3), so there exists a matrix \(A \in \mathbf{GL}(2, k)\) such that \(AxA^{-1} = X\) , and consequently an automorphism \(\alpha\) of \(\mathfrak{sl}(2, k)\) such that \(\alpha(x) = x'\) ; now \(\alpha \in \operatorname{Aut}_e(\mathfrak{g})\) (§5, no. 3, Cor. 2 of Prop. 5). We now pass to the general case; we assume that \(\mathfrak{t}\) and \(\mathfrak{t}'\) are G-conjugate and prove that \(x\) and \(x'\) are G-conjugate. We can assume that \(\mathfrak{t} = \mathfrak{t}'\) . By the preceding, there exists \(\beta \in \operatorname{Aut}_e(\mathfrak{t})\) such that \(\beta x = x'\) . Now, if \(t \in \mathfrak{t}\) is such that \(\operatorname{ad}_{\mathfrak{t}} t\) is nilpotent, then \(\operatorname{ad}_{\mathfrak{g}} t\) is nilpotent; so \(\beta\) extends to an element of \(\operatorname{Aut}_e(\mathfrak{g})\) .

Remark. The three conditions of Prop. 1 are equivalent if we assume only that \(k = k^2\) (cf.~Exerc. 1).

